%% ============================================================
%%  Aula 8 — Aprendizagem por imitação, aprendizado de recompensa e RLHF
%%  Adaptação em pt-BR do CS234 (Stanford, Emma Brunskill), Lecture 8
%%  Prof. Marcelo Keese Albertini — Universidade Federal de Uberlândia
%%  Compilar com XeLaTeX (2x) — latexmk -xelatex aula08.tex
%% ============================================================
\documentclass[aspectratio=169, 11pt]{beamer}
%% .sty do projeto na raiz do repositório (../)
\makeatletter
\def\input@path{{./}{../}}
\makeatother


\usetheme{AprendizadoRL}      % ANTES do babel — fontspec precisa vir primeiro
\usepackage{babel}
\babelprovide[import, main]{brazilian}
\usepackage{booktabs}

\title[Aula 8 — Imitação, recompensa e RLHF]{Aula 8: Aprendizagem por imitação,\\ aprendizado de recompensa e RLHF}
\subtitle[Aprendizagem por Reforço]{Erros compostos \textbullet\ DAGGER \textbullet\ RL inverso \textbullet\ Bradley--Terry \textbullet\ RLHF \textbullet\ DPO}
\author[Prof. Marcelo K. Albertini]{Prof. Marcelo Keese Albertini \\
  \footnotesize Universidade Federal de Uberlândia \\
  \footnotesize adaptado de Emma Brunskill, CS234 (Stanford)}
\institute{Universidade Federal de Uberlândia}
\date{2026}

% ---- Nota discreta: perguntas ficam para o professor conduzir em aula
\newcommand{\paradiscussao}{\smallskip
  {\footnotesize\color{rlCinza}\itshape Pergunta para discussão conduzida em aula.}}

%% ------------------------------------------------------------
%% Macros locais desta aula
%% ------------------------------------------------------------
\newcommand{\thv}{\theta}
\newcommand{\Qhat}{\hat{Q}}
\newcommand{\polth}{\pol_{\thv}}                 % política parametrizada
\newcommand{\polE}{\pol^{E}}                     % política do especialista
\newcommand{\pref}{\pol_{\mathrm{ref}}}          % modelo de referência
\newcommand{\psft}{\pol_{\mathrm{SFT}}}          % modelo após ajuste supervisionado
\newcommand{\rphi}{r_{\phi}}                     % modelo de recompensa aprendido
\newcommand{\Dset}{\mathcal{D}}
\newcommand{\traj}{\tau}
\newcommand{\yw}{y_{w}}                          % resposta preferida
\newcommand{\yl}{y_{\ell}}                       % resposta preterida
\newcommand{\KLd}[2]{\mathrm{KL}\!\left(#1 \,\middle\|\, #2\right)}
\newcommand{\sig}{\sigma}
\newcommand{\Hent}{\mathcal{H}}
\newcommand{\gth}{\nabla_{\thv}}
\newcommand{\dpref}{d^{\pol}}                    % distribuição de estados induzida

\begin{document}

\begin{frame}[plain, noframenumbering]\titlepage\end{frame}

%% ------------------------------------------------------------
\begin{frame}{Onde estamos}{Da recompensa dada à recompensa aprendida}

  \begin{itemize}\setlength{\itemsep}{0.45ex}
    \item \textbf{Aulas 2--3:} com modelo, calcular $\Vstar$ e $\polstar$ por
      programação dinâmica. \textbf{Aula 4:} sem modelo, aprender $\Qhat$ de
      amostras.
    \item \textbf{Aulas 5--6:} parametrizar a política e subir o gradiente do
      retorno --- REINFORCE, ator--crítico, TRPO, PPO.
    \item \textbf{Aula 7:} e se ninguém escrever a recompensa? Primeira resposta:
      \emph{copie um especialista}.
    \item \textbf{Hoje:} por que copiar não basta, como \emph{inferir} a
      recompensa e como usar \emph{preferências humanas} --- o caminho até o
      RLHF e o DPO.
  \end{itemize}

  \begin{ideiabox}
    Até aqui $\rec$ era um dado do problema. A partir de hoje ela vira
    \emph{objeto de aprendizado} --- e quase todos os problemas novos da aula
    nascem exatamente daí.
  \end{ideiabox}
\end{frame}

%% ------------------------------------------------------------
\begin{frame}{Recapitulação da aula passada}

  \begin{quizbox}{}
    Quais afirmações sobre o \textbf{REINFORCE} são verdadeiras?
    \begin{enumerate}\setlength{\itemsep}{0ex}
      \item[(a)] Uma linha de base reduz a variância das atualizações.
      \item[(b)] O método converge para um ótimo \emph{global}.
      \item[(c)] Pode partir de uma política determinística subótima e ainda
        convergir a um ótimo local.
      \item[(d)] Um passo de gradiente pode \emph{piorar} a política.
    \end{enumerate}
    \vspace{-0.3em}\paradiscussao
  \end{quizbox}

  \pause

  \begin{respbox}
    \textbf{(a) e (d).} \;(b): o objetivo é não convexo em $\thv$. \;(c): uma
    política determinística nunca amostra as demais ações, que nunca recebem
    gradiente --- ela fica presa onde começou.
  \end{respbox}
\end{frame}

%% ------------------------------------------------------------
\begin{frame}{Plano de hoje}

  \begin{columns}[T]
  \begin{column}{0.49\textwidth}
    \begin{enumerate}\setlength{\itemsep}{0.5ex}
      \item \textbf{Imitação e erros compostos}\\
        {\small por que clonar comportamento degrada como $T^2$, e o que o
         DAGGER faz a respeito}
      \item \textbf{Aprendizado de recompensa}\\
        {\small RL inverso, a não unicidade de $\rec$, máxima entropia}
      \item \textbf{Humanos no laço}\\
        {\small o espectro de esforço humano e por que comparações são baratas}
    \end{enumerate}
  \end{column}
  \begin{column}{0.49\textwidth}
    \begin{enumerate}\setlength{\itemsep}{0.5ex}
      \setcounter{enumi}{3}
      \item \textbf{Modelo de Bradley--Terry}\\
        {\small de preferências ruidosas a uma recompensa escalar}
      \item \textbf{RLHF}\\
        {\small o \emph{pipeline} de três etapas e a penalidade KL}
      \item \textbf{DPO}\\
        {\small como eliminar o modelo de recompensa por álgebra}
    \end{enumerate}
  \end{column}
  \end{columns}

  \begin{ideiabox}
    O fio condutor: a aula toda responde à mesma pergunta --- ``de onde vem o
    sinal de recompensa quando não há placar?''
  \end{ideiabox}
\end{frame}

%% ------------------------------------------------------------
\begin{frame}{Por que isto virou o assunto central}{RL como etapa final de treinamento de modelos de linguagem}

  \begin{center}
  \scalebox{0.92}{\begin{tikzpicture}[font=\scriptsize, node distance=6mm]
    \node[draw=rlCinza, thick, fill=rlFundo, rounded corners=3pt,
          text width=27mm, align=center, minimum height=13mm] (pt)
      {\textbf{Pré-treino}\\[0.2em] prever o próximo \emph{token}\\ (trilhões de palavras)};
    \node[draw=rlAzulClaro, thick, fill=rlAzulClaro!8, rounded corners=3pt,
          text width=27mm, align=center, minimum height=13mm, right=9mm of pt] (sft)
      {\textbf{Ajuste supervisionado}\\[0.2em] imitar respostas escritas por pessoas};
    \node[draw=rlLaranja, very thick, fill=rlLaranja!10, rounded corners=3pt,
          text width=27mm, align=center, minimum height=13mm, right=9mm of sft] (rl)
      {\textbf{RL sobre preferências}\\[0.2em] otimizar o que as pessoas de fato aprovam};
    \draw[trans] (pt) -- (sft);
    \draw[trans destac] (sft) -- (rl);
    \node[below=3mm of pt, text=rlCinza, text width=27mm, align=center]
      {aprendizado auto-supervisionado};
    \node[below=3mm of sft, text=rlCinza, text width=27mm, align=center]
      {\textbf{clonagem de comportamento}\\ (Aula 7)};
    \node[below=3mm of rl, text=rlLaranja, text width=27mm, align=center]
      {\textbf{hoje}};
  \end{tikzpicture}}
  \end{center}

  \begin{itemize}\setlength{\itemsep}{0.3ex}
    \item As duas primeiras etapas são supervisionadas; a terceira é RL --- e é
      a que exige que exista uma \emph{recompensa}.
    \item Ninguém sabe escrever ``responda de forma útil e honesta'' como uma
      $\rec(s,a)$. Mas todos sabem dizer \emph{qual de duas respostas é melhor}.
  \end{itemize}
\end{frame}

%% ============================================================
\section{Imitação e erros compostos}
%% ============================================================

%% ------------------------------------------------------------
\begin{frame}{O problema de aprendizagem por imitação}

  \begin{defbox}{Aprendizagem por imitação}
    Dados: espaços $\gS$ e $\gA$, possivelmente um modelo de transição
    $\tran{s'}{s,a}$, e um conjunto de demonstrações
    \[
      \Dset = \bigl\{(s_0, a_0, s_1, a_1, \dots)\bigr\},
      \qquad a_t \sim \polE(\cdot \mid s_t),
    \]
    geradas por um \textbf{especialista} $\polE$. \emph{Não há função de
    recompensa.} Objetivo: obter uma política tão boa quanto $\polE$.
  \end{defbox}

  \begin{itemize}\setlength{\itemsep}{0.3ex}
    \item \textbf{Clonagem de comportamento} (Aula 7): trate isso como
      classificação --- aprenda $\pol(a\mid s)$ minimizando a perda
      supervisionada sobre os pares $(s,a)$ do especialista.
    \item Simples, rápido, sem interação com o ambiente --- e responsável por
      resultados históricos (ALVINN, 1989) e atuais (o ``SFT'' dos LLMs).
  \end{itemize}
\end{frame}

%% ------------------------------------------------------------
\begin{frame}{O que a clonagem de comportamento supõe --- e o que ela ignora}

  \begin{columns}[T]
  \begin{column}{0.5\textwidth}
    \begin{defbox}{A hipótese emprestada}
      Aprendizado supervisionado supõe pares $(x,y)$ \textbf{i.i.d.},
      extraídos da \emph{mesma} distribuição no treino e no teste.
    \end{defbox}
    \smallskip
    Em um MDP isso é falso duas vezes: os estados de uma trajetória são
    \textbf{dependentes} no tempo, e a distribuição deles depende da
    \textbf{política que os gera}.
  \end{column}
  \begin{column}{0.48\textwidth}
    \begin{alertabox}
      \textbf{Deslocamento de distribuição.} Treino: $s_t \sim d^{\polE}$;
      teste: $s_t \sim d^{\pol}$. O classificador é bom \emph{onde o
      especialista esteve} --- e um erro basta para sair dessa região.
    \end{alertabox}
  \end{column}
  \end{columns}

  \begin{ideiabox}
    O problema estrutural da imitação, e a razão de ela não ser ``só
    classificação'': as ações do aluno mudam as perguntas da prova.
    {\scriptsize\color{rlCinza}\;(Ross \& Bagnell, 2011.)}
  \end{ideiabox}
\end{frame}

%% ------------------------------------------------------------
\begin{frame}{Erros independentes no tempo: o caso benigno}

  Suponha que, em cada passo, a política clonada erre com probabilidade
  $\varepsilon$, e que \emph{um erro não afete os estados seguintes}:
  \[
    \E[\text{total de erros}] \;\le\; \varepsilon\,T .
  \]

  \begin{center}
  \scalebox{0.95}{\begin{tikzpicture}[font=\scriptsize]
    \draw[trans, rlAzul, very thick] (0,0) -- (9,0);
    \foreach \x in {0,1,...,9}{\fill[rlAzul!35] (\x,0) circle (1.1pt);}
    \foreach \x in {2,6}{
      \fill[rlLaranja] (\x,0) circle (2.2pt);
      \draw[rlLaranja, thick] (\x,0) -- ++(0,-0.45);
      \node[below, text=rlLaranja] at (\x,-0.45) {erro};}
    \node[above, text=rlCinza] at (4.5,0.15) {trajetória do especialista};
    \node[right, text=rlAzul] at (9.1,0) {$t$};
  \end{tikzpicture}}
  \end{center}

  \begin{ideiabox}
    Neste mundo idealizado o agente ``volta ao trilho'' sozinho depois de cada
    erro. O custo é \textbf{linear} no horizonte --- exatamente o que a intuição
    de aprendizado supervisionado prevê.
  \end{ideiabox}

  A pergunta da aula: \emph{o que muda quando um erro desloca a trajetória?}
\end{frame}

%% ------------------------------------------------------------
\begin{frame}{Erros compostos: o custo real}

  \begin{columns}[T]
  \begin{column}{0.53\textwidth}
    Um erro em $t$ leva o agente a um estado que o especialista nunca visitou.
    Ali a política não tem garantia --- e a suposição pessimista é que
    \emph{todos} os $T-t$ passos restantes fiquem comprometidos:
    \begin{align*}
      \E[\text{erros}] &\lesssim
        \varepsilon\bigl(T + (T-1) + \dots + 1\bigr)\\[0.2em]
        &= \varepsilon\,\frac{T(T+1)}{2}
        \;\propto\; \mdestaque{rlLaranja}{\varepsilon\,T^{2}}
    \end{align*}

    \smallskip
    {\scriptsize\color{rlCinza} Argumento heurístico; o resultado formal é o
    Teorema 2.1 de Ross \& Bagnell (AISTATS 2010).}
  \end{column}
  \begin{column}{0.45\textwidth}
    \begin{center}
    \scalebox{0.9}{\begin{tikzpicture}[font=\scriptsize]
      % eixo
      \draw[->, rlCinza] (0,0) -- (4.6,0) node[right] {$T$};
      \draw[->, rlCinza] (0,0) -- (0,3.2) node[above] {erros};
      % linear
      \draw[rlAzulClaro, very thick] (0,0) -- (4.3,1.29);
      \node[text=rlAzulClaro, right] at (3.3,0.75) {$\varepsilon T$};
      % quadratico
      \draw[rlLaranja, very thick] plot[domain=0:4.3, samples=40]
        (\x, {0.16*\x*\x});
      \node[text=rlLaranja, left] at (3.5,2.6) {$\varepsilon T^{2}$};
    \end{tikzpicture}}
    \end{center}
    \begin{alertabox}
      Com $\varepsilon = 1\%$ e $T = 1000$: $10$ contra $\sim\!10^{4}$ erros
      esperados. A política ``com 99\% de acerto'' é inutilizável.
    \end{alertabox}
  \end{column}
  \end{columns}
\end{frame}

%% ------------------------------------------------------------
\begin{frame}{A geometria do problema}

  \begin{center}
  \scalebox{0.95}{\begin{tikzpicture}[font=\scriptsize]
    % faixa de estados visitados pelo especialista
    \fill[rlAzul!7, rounded corners=6pt] (0,-0.55) .. controls (3,-0.15) and (5,0.95) .. (8.4,1.15)
      -- (8.4,1.95) .. controls (5,1.75) and (3,0.65) .. (0,0.25) -- cycle;
    \node[text=rlAzul!70, anchor=west] at (0.15,1.6) {suporte dos dados do especialista};
    % trajetoria do especialista
    \draw[rlAzul, very thick] (0,-0.15) .. controls (3,0.25) and (5,1.35) .. (8.4,1.55);
    \node[text=rlAzul, anchor=west] at (8.5,1.55) {$\polE$};
    % trajetoria do agente
    \draw[rlLaranja, very thick] (0,-0.15) .. controls (2.2,0.1) and (2.9,-0.1) .. (3.4,-0.45)
      .. controls (4.6,-1.3) and (6.2,-2.0) .. (8.4,-2.35);
    \node[text=rlLaranja, anchor=west] at (8.5,-2.35) {$\pol$ clonada};
    % ponto do erro
    \fill[rlVermelho] (3.4,-0.45) circle (2.4pt);
    \node[text=rlVermelho, anchor=north west] at (3.5,-0.55) {primeiro erro};
    % chaves
    \draw[trans, rlCinza] (7.0,-1.9) -- (7.0,0.95);
    \node[text=rlCinza, anchor=west, text width=25mm] at (7.1,-0.5)
      {erro que \emph{cresce}: nenhum dado aqui};
  \end{tikzpicture}}
  \end{center}

  \begin{ideiabox}
    O erro não é ``ruído em torno da resposta certa''. É um \textbf{desvio que se
    realimenta}: quanto mais longe do suporte dos dados, pior a política, e pior
    a política, mais longe ela vai.
  \end{ideiabox}
\end{frame}

%% ------------------------------------------------------------
\begin{frame}{Teste sua compreensão}

  \begin{quizbox}{}
    Uma equipe treina um piloto automático por clonagem de comportamento e o
    carro sai da pista após alguns minutos. A proposta é \emph{coletar dez vezes
    mais demonstrações do motorista}.

    \smallskip
    Isso resolve o problema? Em que sentido ajuda e em que sentido não? Que
    dado atacaria a causa em vez do sintoma?
    \paradiscussao
  \end{quizbox}

  \pause

  \begin{respbox}
    Ajuda a reduzir $\varepsilon$, mas \textbf{não muda o expoente}. Mais dados
    do motorista são mais dados da região \emph{onde ele dirige bem} --- e o
    problema está fora dela. Ataca a causa obter rótulos \textbf{nos estados que
    a política aprendida visita}.
  \end{respbox}
\end{frame}

%% ------------------------------------------------------------
\begin{frame}{DAGGER: agregação de conjuntos de dados}

  \begin{center}
  \scalebox{0.88}{\begin{tikzpicture}[font=\scriptsize, node distance=13mm]
    \node[draw=rlAzul, thick, fill=rlAzul!7, rounded corners=3pt,
          text width=30mm, align=center, minimum height=11mm] (treina)
      {\textbf{Treinar $\pol_i$}\\ por classificação sobre $\Dset$};
    \node[draw=rlLaranja, very thick, fill=rlLaranja!10, rounded corners=3pt,
          text width=30mm, align=center, minimum height=11mm, right=17mm of treina] (roda)
      {\textbf{Executar $\pol_i$}\\ no ambiente, coletar estados visitados};
    \node[draw=rlTeal, thick, fill=rlTeal!10, rounded corners=3pt,
          text width=30mm, align=center, minimum height=11mm, below=11mm of roda] (rotula)
      {\textbf{Perguntar ao especialista}\\ qual ação ele tomaria \emph{ali}};
    \node[draw=rlAzulClaro, thick, fill=rlAzulClaro!8, rounded corners=3pt,
          text width=30mm, align=center, minimum height=11mm, below=11mm of treina] (agrega)
      {\textbf{Agregar}\\ $\Dset \leftarrow \Dset \cup \Dset_i$};
    \draw[trans destac] (treina) -- (roda);
    \draw[trans destac] (roda) -- (rotula);
    \draw[trans destac] (rotula) -- (agrega);
    \draw[trans destac] (agrega) -- (treina);
  \end{tikzpicture}}
  \end{center}

  \begin{itemize}\setlength{\itemsep}{0.3ex}
    \item A ideia: \emph{obter rótulos do especialista ao longo do caminho que a
      política aprendida percorre}.
    \item Resultado: uma política determinística com bom desempenho \textbf{sob
      a distribuição de estados que ela mesma induz}.
    \item É uma redução a aprendizado \emph{online} sem arrependimento --- o
      custo volta a ser $O(\varepsilon T)$.
  \end{itemize}
\end{frame}

%% ------------------------------------------------------------
\begin{frame}{DAGGER: o algoritmo}

  \begin{algbox}{DAGGER (Ross et al., 2011)}
    {\footnotesize
    \begin{itemize}\setlength{\itemsep}{0.15ex}
      \item Inicialize $\Dset \leftarrow \emptyset$ e $\pol_1$ (por exemplo, o
        próprio $\polE$ ou uma clonagem inicial)
      \item \textbf{para} $i = 1, 2, \dots, N$ \textbf{faça}
        \begin{itemize}\setlength{\itemsep}{0.1ex}
          \item Defina $\hat\pol_i = \beta_i\,\polE + (1-\beta_i)\,\pol_i$
            \hfill{\scriptsize(mistura; $\beta_i \to 0$)}
          \item Execute $\hat\pol_i$ e colete as trajetórias de estados visitados
          \item Consulte o especialista em cada estado visitado:
            $\Dset_i = \{(s, \polE(s))\}$
          \item Agregue: $\Dset \leftarrow \Dset \cup \Dset_i$
          \item Treine $\pol_{i+1}$ minimizando a perda de classificação sobre
            \emph{todo} o $\Dset$
        \end{itemize}
      \item \textbf{fim para}: devolva a melhor $\pol_i$ segundo validação
    \end{itemize}}
  \end{algbox}

  \begin{ideiabox}
    A mistura $\beta_i$ evita que as primeiras iterações passem tempo demais em
    estados catastróficos, quando $\pol_i$ ainda é ruim --- a lógica do
    instrutor de direção que segura o volante no início.
  \end{ideiabox}
\end{frame}

%% ------------------------------------------------------------
\begin{frame}{DAGGER: a limitação decisiva}

  \begin{quizbox}{}
    O DAGGER converte um erro $O(\varepsilon T^2)$ em $O(\varepsilon T)$. Qual é
    o preço --- e por que ele inviabiliza o método em boa parte das aplicações
    modernas?
    \paradiscussao
  \end{quizbox}

  \pause

  \begin{respbox}
    Ele exige um \textbf{especialista disponível durante todo o treinamento},
    capaz de rotular estados arbitrários sob demanda.
    \begin{itemize}\setlength{\itemsep}{0.1ex}
      \item \textbf{Custo}: a supervisão escala com estados visitados, não com
        demonstrações.
      \item \textbf{Contexto}: perguntar ``o que você faria \emph{agora}?'' num
        quadro congelado produz rótulos ruins.
      \item \textbf{Segurança}: para rotular um estado é preciso chegar até ele
        --- e alguns não devem ser visitados.
    \end{itemize}
  \end{respbox}
\end{frame}

%% ------------------------------------------------------------
\begin{frame}{Balanço da imitação}

  \begin{center}\footnotesize
  \begin{tabular}{@{}p{0.27\linewidth}p{0.2\linewidth}p{0.2\linewidth}p{0.2\linewidth}@{}}
    \toprule
    & \textbf{Clonagem} & \textbf{DAGGER} & \textbf{RL inverso} \\
    \midrule
    Erro no horizonte & $O(\varepsilon T^2)$ & $O(\varepsilon T)$ &
      depende do planejador \\
    Interage com o ambiente & não & sim & sim \\
    Precisa do especialista \emph{durante} o treino & não & \textbf{sim} & não \\
    O que produz & uma política & uma política & uma \textbf{recompensa} \\
    Generaliza para novas dinâmicas & não & não & \textbf{sim} \\
    \bottomrule
  \end{tabular}
  \end{center}

  \begin{ideiabox}
    A última linha é o argumento a favor de aprender a \emph{recompensa} em vez
    da política: ela é a descrição mais \textbf{compacta e transferível} da
    tarefa. Mude o carro ou a estrada e a política precisa ser refeita --- a
    intenção, não.
  \end{ideiabox}
\end{frame}

%% ============================================================
\section{Aprendizado de recompensa}
%% ============================================================

%% ------------------------------------------------------------
\begin{frame}{Inverter o problema}

  \begin{defbox}{Aprendizado por reforço inverso}
    Dados $\gS$, $\gA$, o modelo de transição $\tran{s'}{s,a}$ e demonstrações
    $(s_0, a_0, s_1, a_1, \dots)$ com $a_t \sim \polE(\cdot\mid s_t)$, mas
    \textbf{sem} função de recompensa: inferir uma $\rec$ sob a qual o
    comportamento observado seja ótimo.
  \end{defbox}

  \begin{center}
  \scalebox{0.95}{\begin{tikzpicture}[font=\scriptsize, node distance=8mm]
    \node[draw=rlAzul, thick, fill=rlAzul!7, rounded corners=3pt,
          text width=24mm, align=center, minimum height=9mm] (r) {$\rec$};
    \node[draw=rlAzul, thick, fill=rlAzul!7, rounded corners=3pt,
          text width=24mm, align=center, minimum height=9mm, right=22mm of r] (p) {$\polstar$};
    \draw[trans] (r) -- node[above, text=rlCinza] {RL \emph{direto}: planejar} (p);
    \draw[trans destac] ([yshift=-5mm]p.south) -- node[below, text=rlLaranja]
      {RL \emph{inverso}: explicar} ([yshift=-5mm]r.south);
    \draw[rlLaranja, thick] (p.south) -- ([yshift=-5mm]p.south);
    \draw[rlLaranja, thick] (r.south) -- ([yshift=-5mm]r.south);
  \end{tikzpicture}}
  \end{center}

  \begin{itemize}\setlength{\itemsep}{0.25ex}
    \item A recompensa é a descrição \textbf{mais sucinta} da tarefa (Ng \&
      Russell, 2000). Uma vez recuperada, pode-se replanejar em outra dinâmica.
    \item Antes de qualquer algoritmo, uma pergunta de identificabilidade.
  \end{itemize}
\end{frame}

%% ------------------------------------------------------------
\begin{frame}{Teste sua compreensão}

  \begin{quizbox}{}
    Suponha que a política do especialista $\polE$ seja \emph{ótima}. Sobre a
    recompensa $\rec$ que a explica:
    \begin{enumerate}\setlength{\itemsep}{0.1ex}
      \item[1.] existe uma única $\rec$ que torna $\polE$ ótima;
      \item[2.] existem muitas $\rec$ que tornam $\polE$ ótima;
      \item[3.] depende do MDP.
    \end{enumerate}
    Justifique exibindo, se possível, um contraexemplo.
    \paradiscussao
  \end{quizbox}

  \pause

  \begin{respbox}
    \textbf{Existe um conjunto infinito de $\rec$ compatíveis.} O caso mais
    brutal: $\rec \equiv 0$ torna \emph{toda} política ótima. Sem restrições, o
    problema é mal posto --- e todo o RL inverso consiste em escolher
    \emph{qual} solução preferir.
  \end{respbox}
\end{frame}

%% ------------------------------------------------------------
\begin{frame}{As três fontes de ambiguidade}

  \begin{columns}[T]
  \begin{column}{0.5\textwidth}
    \begin{alertabox}
      \textbf{1. Degenerescência.} $\rec\equiv 0$ (ou qualquer constante)
      explica qualquer comportamento.
    \end{alertabox}
    \begin{alertabox}
      \textbf{2. Escala.} Se $\rec$ explica $\polE$, então $c\,\rec$ com
      $c>0$ também: $\argmax$ não muda.
    \end{alertabox}
  \end{column}
  \begin{column}{0.48\textwidth}
    \begin{alertabox}
      \textbf{3. Modelagem por potencial.} Para qualquer $\Phi:\gS\to\R$,
      \[
        \rec'(s,a,s') = \rec(s,a,s') + \gamma\Phi(s') - \Phi(s)
      \]
      preserva \emph{exatamente} o conjunto de políticas ótimas
      {\scriptsize(Ng, Harada \& Russell, 1999)}.
    \end{alertabox}
  \end{column}
  \end{columns}

  \begin{ideiabox}
    Guarde a invariância 2: ela reaparece no modelo de Bradley--Terry e é a razão
    de o modelo de recompensa de um LLM não ter escala absoluta interpretável.
    Como escolher entre infinitas soluções? Impondo uma \textbf{estrutura} e um
    \textbf{critério de desempate}.
  \end{ideiabox}
\end{frame}

%% ------------------------------------------------------------
\begin{frame}{Estrutura: recompensa linear em atributos}

  Suponha $\rec(s) = w^{\!\top}\!x(s)$, com atributos $x(s)\in\R^{n}$ conhecidos.
  O valor vira o produto interno com a \textbf{esperança de atributos}:
  \[
    \Vpi = \E\Bigl[\sum_{t\ge 0}\gamma^{t}\, w^{\!\top} x(s_t) \,\Bigm|\, \pol\Bigr]
      = w^{\!\top} \underbrace{\E\Bigl[\sum_{t\ge 0}\gamma^{t} x(s_t) \Bigm| \pol\Bigr]}_{\textstyle \mu(\pol)}
      = w^{\!\top}\mu(\pol).
  \]

  \begin{teobox}{Consequência}
    Se $\norm{w}_1 \le 1$ e $\norm{\mu(\pol) - \mu(\polE)}_1 \le \epsilon$, então
    $\abs{ w^{\!\top}\mu(\pol) - w^{\!\top}\mu(\polE)} \le \epsilon$ para todo
    $w$ admissível: \textbf{casar as esperanças de atributos basta}.
  \end{teobox}

  \begin{ideiabox}
    Em vez de descobrir ``o que o especialista queria'', basta
    \emph{comportar-se como ele estatisticamente}.
  \end{ideiabox}
\end{frame}

%% ------------------------------------------------------------
\begin{frame}{Desempate: o princípio da máxima entropia}

  Muitas distribuições sobre trajetórias casam as esperanças de atributos.
  \textbf{Escolha a menos comprometida com o que não foi observado} --- a de
  maior entropia:
  \[
    \max_{P}\; -\sum_{\traj} P(\traj)\log P(\traj)
    \quad\text{sujeito a}\quad
    \sum_{\traj} P(\traj)\, x(\traj) = \tilde{x}_{E},
    \quad \sum_{\traj}P(\traj) = 1 .
  \]

  \pause
  A solução por multiplicadores de Lagrange é uma família exponencial:
  \[
    \mdestaque{rlLaranja}{P(\traj \mid w) \;=\; \frac{1}{Z(w)}\,
      \exp\bigl(w^{\!\top} x(\traj)\bigr)},
    \qquad
    x(\traj) = \sum_{t} \gamma^{t} x(s_t).
  \]

  \begin{ideiabox}
    Leitura: o especialista é \textbf{aproximadamente} ótimo. Trajetórias
    melhores são exponencialmente mais prováveis, mas erros não são impossíveis
    --- o que descreve muito melhor uma pessoa do que a hipótese de otimalidade
    exata.
  \end{ideiabox}

  {\scriptsize\color{rlCinza}Ziebart et al., \emph{Maximum Entropy Inverse
  Reinforcement Learning}, AAAI 2008.}
\end{frame}

%% ------------------------------------------------------------
\begin{frame}{MaxEnt IRL: como se aprende $w$}

  Máxima verossimilhança sobre as demonstrações $\Dset$, com
  $w^{*} = \argmax_{w} \sum_{\traj \in \Dset} \log P(\traj\mid w)$:
  \[
    \nabla_w \mathcal{L} = \underbrace{\tilde{x}_{E}}_{\text{atributos observados}}
      - \underbrace{\E_{\traj \sim P(\cdot\mid w)}\bigl[x(\traj)\bigr]}_{\text{atributos do modelo atual}} .
  \]

  \begin{algbox}{MaxEnt IRL}
    {\footnotesize
    \textbf{Repita} até convergir: (i) \emph{laço interno} --- resolva o MDP com
    a recompensa atual $w^{\!\top}x(s)$ e obtenha $\pol$ e as visitas esperadas
    de estados; (ii) calcule o gradiente acima e atualize $w$.}
  \end{algbox}

  \begin{alertabox}
    O laço interno é \emph{resolver um MDP inteiro a cada passo de gradiente}.
    Guided Cost Learning (Finn et al., 2016) o troca por amostragem por
    importância --- e revela a equivalência com GANs: a recompensa é o
    discriminador.
  \end{alertabox}
\end{frame}

%% ------------------------------------------------------------
\begin{frame}{Balanço: aprendizagem por imitação}

  \begin{itemize}\setlength{\itemsep}{0.4ex}
    \item \textbf{É poderosa.} Demonstrações são a forma mais natural de
      comunicar uma tarefa a uma máquina, e existem em abundância.
    \item \textbf{Tem muitas extensões}: imitação a partir de observação (sem
      ações), imitação com um único exemplo, imitação a partir de vídeo humano.
    \item \textbf{Máxima entropia é a ideia central} a reter: transformar
      ``explicar comportamento'' em um problema de verossimilhança bem posto.
    \item \textbf{Mas} pressupõe um especialista que \emph{sabe fazer} a tarefa.
      E se ninguém souber --- se as pessoas só souberem \emph{reconhecer} uma
      execução boa?
  \end{itemize}

  \begin{ideiabox}
    Essa última pergunta é a ponte para o resto da aula. Escrever um soneto é
    difícil; dizer qual de dois sonetos é melhor, nem tanto. O RLHF é construído
    inteiramente sobre essa assimetria.
  \end{ideiabox}
\end{frame}

%% ============================================================
\section{Pessoas no laço de treinamento}
%% ============================================================

%% ------------------------------------------------------------
\begin{frame}{Como uma pessoa pode ajudar a treinar um agente}

  \begin{center}
  \scalebox{0.9}{\begin{tikzpicture}[font=\scriptsize]
    \draw[trans, very thick, rlCinza] (0,0) -- (11,0);
    \node[below, text=rlCinza] at (5.5,-0.9) {\textbf{esforço humano}};
    \node[right, text=rlCinza] at (11.05,0) {alto};
    \node[left, text=rlCinza] at (-0.05,0) {baixo};

    \foreach \x/\rot/\cor in {
      1.6/{Demonstrações\\ (uma vez)}/rlAzulClaro,
      5.5/{Rótulos par a par\\ (contínuos, baratos)}/rlLaranja,
      9.4/{DAGGER / ensino\\ constante}/rlVermelho}{
      \fill[\cor] (\x,0) circle (3pt);
      \draw[\cor, thick] (\x,0) -- ++(0,0.55);
      \node[above, text=\cor, align=center, text width=28mm] at (\x,0.55) {\rot};}

    \node[below, text width=30mm, align=center, text=rlCinza] at (1.6,-0.15)
      {só funciona no suporte dos dados};
    \node[below, text width=30mm, align=center, text=rlLaranja] at (5.5,-0.15)
      {\textbf{ponto de equilíbrio?}};
    \node[below, text width=30mm, align=center, text=rlCinza] at (9.4,-0.15)
      {não escala};
  \end{tikzpicture}}
  \end{center}

  \begin{itemize}\setlength{\itemsep}{0.25ex}
    \item Há muitas formas de uma pessoa treinar um agente: demonstrar, corrigir,
      pontuar, comparar, restringir, interromper.
    \item A escolha importa quando queremos agentes que igualem o desempenho
      humano \textbf{e} respeitem valores humanos.
  \end{itemize}
\end{frame}

%% ------------------------------------------------------------
\begin{frame}{Recompensa dada por pessoas: o caminho ingênuo}

  \begin{itemize}\setlength{\itemsep}{0.35ex}
    \item \textbf{TAMER} (Knox \& Stone, 2008): a pessoa aperta ``bom'' ou
      ``ruim'' enquanto o agente age; treina-se um modelo do reforço humano.
    \item \textbf{Robôs ensináveis} (Thomaz \& Breazeal, 2008): pessoas
      \emph{não} dão reforço como a teoria supõe --- elas antecipam, premiam
      intenção e usam o canal para \emph{orientar}, não só para avaliar.
  \end{itemize}

  \begin{alertabox}
    Pedir números diretamente falha por três razões: \textbf{calibração} (o
    ``7'' de uma pessoa não é o de outra, nem o dela na semana seguinte),
    \textbf{ruído} (julgamentos absolutos têm variância alta) e \textbf{esforço}
    (exige atenção contínua).
  \end{alertabox}

  \begin{ideiabox}
    A saída, usada em psicometria muito antes de existir RL: não pergunte
    \emph{quanto}; pergunte \emph{qual dos dois}.
  \end{ideiabox}
\end{frame}

%% ------------------------------------------------------------
\begin{frame}{Comparações par a par}

  \vspace{-0.2em}
  \begin{columns}[T]
  \begin{column}{0.49\textwidth}
    \begin{defbox}{Por que comparar é mais fácil}
      \begin{itemize}\setlength{\itemsep}{0.15ex}
        \item Mais fácil que escrever uma recompensa à mão.
        \item Mais confiável que atribuir nota escalar.
        \item Exige só julgamento \emph{relativo}, estável entre avaliadores.
      \end{itemize}
    \end{defbox}
  \end{column}
  \begin{column}{0.49\textwidth}
    \begin{itemize}\setlength{\itemsep}{0.3ex}
      \item Recomendação usa isso há décadas: comparar duas ordenações por
        \emph{interleaving} é mais sensível que medir cada uma.
      \item Pode-se \textbf{escolher quais pares perguntar}: aprendizado ativo
        de preferências {\scriptsize(Sadigh et al., 2017)}.
    \end{itemize}
  \end{column}
  \end{columns}

  \begin{alertabox}
    O que se perde: comparações dão apenas uma \textbf{ordem}, e RL precisa de um
    \emph{número}. Reconstruir a escala a partir da ordem é o papel do modelo de
    Bradley--Terry.
  \end{alertabox}
\end{frame}

%% ============================================================
\section{Do par ao número: Bradley--Terry}
%% ============================================================

%% ------------------------------------------------------------
\begin{frame}{O modelo de Bradley--Terry (1952)}

  Comece pelo caso mais simples: $k$ ações $b_1,\dots,b_k$, sem estado nem
  contexto (um \emph{bandido} de $k$ braços). Suponha que a pessoa faça
  comparações \textbf{ruidosas} e que exista uma recompensa latente
  $r(\cdot)$ com
  \[
    \mdestaque{rlLaranja}{\Prob(b_i \succ b_j)
      = \frac{\exp\bigl(r(b_i)\bigr)}{\exp\bigl(r(b_i)\bigr) + \exp\bigl(r(b_j)\bigr)}
      = \sig\bigl(r(b_i) - r(b_j)\bigr) = p_{ij}}
  \]

  \begin{columns}[T]
  \begin{column}{0.52\textwidth}
    \begin{itemize}\setlength{\itemsep}{0.2ex}
      \item É uma \emph{softmax} entre dois itens --- ou, de forma equivalente,
        uma regressão logística sobre a \textbf{diferença} de recompensas.
      \item \textbf{Transitividade estocástica}: $p_{ik}$ fica determinado por
        $p_{ij}$ e $p_{jk}$. Uma hipótese forte, e testável.
    \end{itemize}
  \end{column}
  \begin{column}{0.46\textwidth}
    \begin{center}
    \scalebox{0.85}{\begin{tikzpicture}[font=\scriptsize]
      \draw[->, rlCinza] (-2.4,0) -- (2.6,0) node[right] {$r(b_i)-r(b_j)$};
      \draw[->, rlCinza] (0,-0.15) -- (0,1.9) node[above] {$p_{ij}$};
      \draw[rlCinza, dashed] (-2.4,1.6) -- (2.4,1.6) node[right, text=rlCinza] {$1$};
      \draw[rlCinza, dashed] (-2.4,0.8) -- (2.4,0.8);
      \node[left, text=rlCinza] at (0,0.8) {$\tfrac12$};
      \draw[rlLaranja, very thick] plot[domain=-2.4:2.4, samples=60]
        (\x, {1.6/(1+exp(-2.2*\x))});
    \end{tikzpicture}}
    \end{center}
  \end{column}
  \end{columns}
\end{frame}

%% ------------------------------------------------------------
\begin{frame}{A recompensa latente não é única --- de novo}

  \begin{teobox}{Invariância a translação}
    Para qualquer constante $c$, substituir $r(\cdot)$ por $r(\cdot) + c$ deixa
    \emph{todas} as probabilidades $p_{ij}$ inalteradas, pois só entra a
    diferença $r(b_i) - r(b_j)$.
  \end{teobox}

  \begin{itemize}\setlength{\itemsep}{0.3ex}
    \item Já sabíamos: sem hipóteses adicionais, o modelo latente de recompensa
      não é identificável (mesma lição do RL inverso).
    \item Bradley--Terry é justamente a \textbf{hipótese estrutural} que reduz a
      indeterminação a uma única constante --- normalizável fixando, por
      exemplo, $\sum_i r(b_i) = 0$.
  \end{itemize}

  \begin{alertabox}
    Consequência prática: o valor absoluto que um modelo de recompensa de LLM
    devolve \textbf{não significa nada}. Só diferenças entre respostas ao
    \emph{mesmo} prompt são interpretáveis --- e por isso a normalização por
    prompt é padrão na implementação.
  \end{alertabox}
\end{frame}

%% ------------------------------------------------------------
\begin{frame}{Que item é ``o melhor''? Três respostas}

  \begin{defbox}{Vencedor de Condorcet}
    $b_i$ tal que $\Prob(b_i \succ b_j) > 0{,}5$ para \emph{todo} $j \ne i$.
    Pode não existir.
  \end{defbox}

  \begin{columns}[T]
  \begin{column}{0.49\textwidth}
    \begin{defbox}{Vencedor de Copeland}
      $b_i$ com o maior número de vitórias par a par: pontuação igual ao número
      de itens que ele vence menos o número dos que o vencem.
    \end{defbox}
  \end{column}
  \begin{column}{0.49\textwidth}
    \begin{defbox}{Vencedor de Borda}
      $b_i$ que maximiza a pontuação esperada, onde a pontuação contra $b_j$ é
      $1$ se $b_i \succ b_j$, $0{,}5$ em caso de empate, e $0$ caso contrário.
    \end{defbox}
  \end{column}
  \end{columns}

  \begin{ideiabox}
    Os algoritmos de bandidos ``em duelo'' historicamente buscavam o vencedor de
    Copeland. \textbf{Sob Bradley--Terry os três critérios coincidem}, pois a
    transitividade estocástica elimina ciclos --- é quando o modelo \emph{falha}
    que a escolha importa. {\scriptsize\color{rlCinza}(Yue et al., JCSS 2012.)}
  \end{ideiabox}
\end{frame}

%% ------------------------------------------------------------
\begin{frame}{Ajustando os parâmetros: entropia cruzada}

  Suponha $N$ tuplas $(b_i, b_j, \mu)$, com $\mu = 1$ se a pessoa marcou
  $b_i \succ b_j$, $\mu = 0{,}5$ se houve empate e $\mu = 0$ caso contrário.
  Maximizar a verossimilhança equivale a minimizar a entropia cruzada:
  \[
    \mathcal{L}(r) = -\!\!\sum_{(b_i,b_j,\mu)\in\Dset}\!\!
      \Bigl[\mu \log \Prob(b_i \succ b_j) + (1-\mu)\log \Prob(b_j \succ b_i)\Bigr].
  \]

  \begin{ideiabox}
    Note o que isto é: \textbf{regressão logística} com atributos $e_i - e_j$ e
    coeficientes $r$ --- convexa, com solução única dada a normalização.
  \end{ideiabox}

  \begin{alertabox}
    Se os dados forem \emph{separáveis} (um item nunca perde), a verossimilhança
    empurra $r \to \infty$: regularização não é opcional.
  \end{alertabox}
\end{frame}

%% ------------------------------------------------------------
\begin{frame}{De itens para trajetórias}

  A mesma construção vale trocando ``item'' por \emph{trajetória}. Sejam
  $\traj^1 = (s_0, a_7, s_{14},\dots)$ e $\traj^2 = (s_0, a_6, s_{12},\dots)$, com
  somas de recompensas latentes $R^1 = \sum_{i} r^1_i$ e $R^2 = \sum_i r^2_i$:
  \[
    \hat{\Prob}\bigl(\traj^1 \succ \traj^2\bigr)
      = \frac{\exp\bigl(\sum_i r^1_i\bigr)}
             {\exp\bigl(\sum_i r^1_i\bigr) + \exp\bigl(\sum_i r^2_i\bigr)}
      = \sig\bigl(R^1 - R^2\bigr).
  \]

  \begin{itemize}\setlength{\itemsep}{0.3ex}
    \item Agora $r$ é uma \textbf{rede neural} $\rphi(s,a)$, treinada pela mesma
      perda de entropia cruzada.
    \item Com $\rphi$ em mãos, o problema volta a ser RL comum: rode PPO com a
      recompensa aprendida.
  \end{itemize}

  \begin{alertabox}
    Repare no que a soma implica: a preferência é atribuída à trajetória
    \emph{inteira}, mas o modelo aprende uma recompensa \emph{por passo}. Este é
    um problema de \textbf{atribuição de crédito} --- o mesmo que já vimos, agora
    dentro do sinal de supervisão.
  \end{alertabox}
\end{frame}

%% ------------------------------------------------------------
\begin{frame}{O resultado que abriu a porta}

  \begin{columns}[T]
  \begin{column}{0.55\textwidth}
    \textbf{Christiano et al. (2017)}, \emph{Deep RL from Human Preferences}:
    \begin{itemize}\setlength{\itemsep}{0.3ex}
      \item Tarefa: fazer um agente simulado executar um \emph{backflip} ---
        comportamento para o qual ninguém sabe escrever uma recompensa.
      \item Interface: mostrar dois clipes curtos, perguntar qual está mais
        próximo do desejado.
      \item Custo total: cerca de \textbf{900 comparações} de um avaliador
        humano.
      \item O modelo de recompensa é treinado \emph{em paralelo} com a política,
        e as consultas se concentram onde ele está mais incerto.
    \end{itemize}
  \end{column}
  \begin{column}{0.43\textwidth}
    \begin{ideiabox}
      Novecentos bits de feedback humano compraram um comportamento que nenhuma
      função de recompensa escrita à mão havia produzido. É a demonstração de
      que \emph{avaliar} pode ser muito mais barato que \emph{especificar}.
    \end{ideiabox}
    \begin{alertabox}
      E é o mesmo algoritmo, sem mudança conceitual, que cinco anos depois
      alinharia modelos de linguagem.
    \end{alertabox}
  \end{column}
  \end{columns}
\end{frame}

%% ============================================================
\section{RLHF: preferências em escala}
%% ============================================================

%% ------------------------------------------------------------
\begin{frame}{O \emph{pipeline} de três etapas}

  \begin{center}
  \scalebox{0.86}{\begin{tikzpicture}[font=\scriptsize, node distance=8mm]
    \node[draw=rlAzulClaro, thick, fill=rlAzulClaro!8, rounded corners=3pt,
          text width=32mm, align=center, minimum height=17mm] (a)
      {\textbf{1. Ajuste por instruções}\\[0.3em]
       clonagem de comportamento sobre demonstrações escritas\\[0.2em]
       $\Rightarrow\ \psft$};
    \node[draw=rlTeal, thick, fill=rlTeal!10, rounded corners=3pt,
          text width=32mm, align=center, minimum height=17mm, right=11mm of a] (b)
      {\textbf{2. Modelo de recompensa}\\[0.3em]
       Bradley--Terry sobre pares $(\yw,\yl)$ rotulados por pessoas\\[0.2em]
       $\Rightarrow\ \rphi$};
    \node[draw=rlLaranja, very thick, fill=rlLaranja!10, rounded corners=3pt,
          text width=32mm, align=center, minimum height=17mm, right=11mm of b] (c)
      {\textbf{3. Otimização por RL}\\[0.3em]
       PPO maximizando $\rphi$ com penalidade KL contra $\pref$\\[0.2em]
       $\Rightarrow\ \polth$};
    \draw[trans destac] (a) -- (b);
    \draw[trans destac] (b) -- (c);
    \draw[trans, dashed] (a.south) .. controls +(0,-1.1) and +(0,-1.1) ..
      node[below, text=rlCinza, pos=0.5] {$\pref \defeq \psft$: ponto de partida \emph{e} âncora} (c.south);
  \end{tikzpicture}}
  \end{center}

  \begin{itemize}\setlength{\itemsep}{0.25ex}
    \item A etapa 1 é exatamente a clonagem de comportamento da Aula 7.
    \item As etapas 2 e 3 são as duas metades da pergunta ``como maximizar algo
      que não sabemos escrever?''
  \end{itemize}
\end{frame}

%% ------------------------------------------------------------
\begin{frame}{Etapa 2: aprender o modelo de recompensa}

  Coleta: para um prompt $x$, gerar várias respostas com $\psft$ e pedir que uma
  pessoa ordene um par. Notação: $\yw$ preferida (\emph{winning}), $\yl$
  preterida (\emph{losing}).

  \[
    \mathcal{L}(\phi) = -\,\E_{(x,\yw,\yl)\sim\Dset}
      \Bigl[\log \sig\bigl(\rphi(x,\yw) - \rphi(x,\yl)\bigr)\Bigr]
  \]

  \begin{columns}[T]
  \begin{column}{0.52\textwidth}
    \begin{itemize}\setlength{\itemsep}{0.2ex}
      \item $\rphi$ costuma ser o próprio LLM com a camada final trocada por uma
        saída escalar.
      \item Uma resposta inteira recebe \textbf{um} número: a soma de recompensas
        por token do slide anterior colapsa numa recompensa terminal.
    \end{itemize}
  \end{column}
  \begin{column}{0.46\textwidth}
    \begin{ideiabox}
      \textbf{Valide o modelo de recompensa antes de otimizar contra ele.}
      Métrica: acurácia em prever julgamentos humanos retidos. Modelos grandes
      treinados com dados suficientes se aproximam da concordância entre dois
      humanos --- que é o teto real.
    \end{ideiabox}
  \end{column}
  \end{columns}

  {\scriptsize\color{rlCinza}Stiennon et al., \emph{Learning to summarize from
  human feedback}, NeurIPS 2020.}
\end{frame}

%% ------------------------------------------------------------
\begin{frame}{Etapa 3: otimizar sem destruir o modelo}

  \[
    \max_{\polth}\;
      \E_{x\sim\Dset,\; y\sim\polth(\cdot\mid x)}
      \Bigl[\, \termo{\rphi(x,y)}{recompensa aprendida}
        \;-\; \beta\, \termo{\log \frac{\polth(y\mid x)}{\pref(y\mid x)}}{
          penalidade por se afastar} \Bigr]
  \]

  \begin{itemize}\setlength{\itemsep}{0.3ex}
    \item Em esperança o segundo termo é
      $\beta\,\KLd{\polth(\cdot\mid x)}{\pref(\cdot\mid x)}$; paga-se um preço
      sempre que $\polth(y\mid x) > \pref(y\mid x)$.
    \item Otimização: PPO, tratando cada token como uma ação.
  \end{itemize}

  \begin{alertabox}
    \textbf{Por que a penalidade é indispensável.} $\rphi$ é um modelo ajustado a
    dados finitos, válido apenas na região onde $\psft$ gera texto. Maximizá-la
    sem restrição leva o modelo a \emph{explorar as falhas do avaliador} ---
    o fenômeno de \emph{reward hacking}. A KL mantém a política dentro do
    domínio onde $\rphi$ ainda é confiável.
  \end{alertabox}
\end{frame}

%% ------------------------------------------------------------
\begin{frame}{Sobre-otimização: a lei de Goodhart, com números}

  \begin{columns}[T]
  \begin{column}{0.47\textwidth}
    \begin{center}
    \scalebox{0.9}{\begin{tikzpicture}[font=\scriptsize]
      \draw[->, rlCinza] (0,0) -- (5.0,0);
      \node[text=rlCinza, below] at (2.5,-0.15) {$\sqrt{\KLd{\polth}{\pref}}$};
      \draw[->, rlCinza] (0,0) -- (0,3.1) node[above, align=center] {qualidade};
      % proxy: cresce sempre
      \draw[rlAzulClaro, very thick] plot[domain=0:4.2, samples=50]
        (\x, {2.7*(1-exp(-0.7*\x))});
      \node[text=rlAzulClaro, above right, font=\tiny] at (1.9,2.35) {recompensa $\rphi$ (medida)};
      % real: sobe e desce
      \draw[rlLaranja, very thick] plot[domain=0:3.8, samples=50]
        (\x, {2.4*(1-exp(-0.9*\x)) - 0.16*\x*\x});
      \node[text=rlLaranja, below, font=\tiny] at (3.0,0.55) {qualidade real};
      \draw[rlVermelho, dashed, thick] (1.9,0) -- (1.9,2.1);
      \node[text=rlVermelho, above, align=center, font=\scriptsize] at (1.9,2.15) {ponto de\\ virada};
    \end{tikzpicture}}
    \end{center}
  \end{column}
  \begin{column}{0.51\textwidth}
    \begin{ideiabox}
      ``Quando uma medida se torna um alvo, ela deixa de ser uma boa medida.''
      $\rphi$ é um \textbf{substituto} da preferência humana; otimizá-lo demais
      melhora o substituto e piora o objeto.
    \end{ideiabox}
    \begin{itemize}\setlength{\itemsep}{0.2ex}
      \item O coeficiente $\beta$ é, na prática, o controle desse compromisso.
      \item A distância KL entre política e referência é a variável natural do
        eixo horizontal --- não o número de passos de otimização.
    \end{itemize}
  \end{column}
  \end{columns}
\end{frame}

%% ------------------------------------------------------------
\begin{frame}{O que se observou na prática}

  \begin{itemize}\setlength{\itemsep}{0.4ex}
    \item \textbf{RLHF supera pré-treino $+$ ajuste supervisionado} em avaliação
      humana --- e a vantagem cresce com a escala do modelo
      {\scriptsize(Stiennon et al., 2020)}.
    \item \textbf{InstructGPT} (Ouyang et al., 2022) levou o \emph{pipeline} a
      dezenas de milhares de tarefas; um modelo muito menor, alinhado, foi
      preferido a um modelo bem maior sem alinhamento.
    \item \textbf{Comparações controladas} (Dubois et al., 2023), usando GPT-4
      como avaliador substituto: PPO funciona, mas linhas de base simples ---
      \emph{best-of-$n$} e treinar sobre as saídas ``boas'' --- vão
      surpreendentemente longe.
  \end{itemize}

  \begin{ideiabox}
    A lição de engenharia: boa parte do ganho do RLHF vem de \textbf{ter um bom
    modelo de recompensa}, não da sofisticação do otimizador. Isso levanta uma
    pergunta natural --- se o RL é a parte cara e frágil, dá para eliminá-la?
  \end{ideiabox}
\end{frame}

%% ============================================================
\section{DPO: preferência direta}
%% ============================================================

%% ------------------------------------------------------------
\begin{frame}{O que incomoda no RLHF}

  \begin{columns}[T]
  \begin{column}{0.49\textwidth}
    \begin{alertabox}
      \begin{itemize}\setlength{\itemsep}{0ex}
        \item \textbf{Quatro modelos na memória}: política, referência,
          recompensa e crítico.
        \item \textbf{Amostragem online} a cada passo --- lento e caro.
        \item \textbf{PPO é sensível} a passo, recorte e normalização.
        \item Duas etapas de treino, com erros que se compõem.
        \vspace{-0.6em}
      \end{itemize}
    \end{alertabox}
  \end{column}
  \begin{column}{0.49\textwidth}
    \begin{ideiabox}
      O modelo de recompensa é só um \emph{intermediário} entre preferências e
      política. Se houver relação em forma fechada entre recompensa e política
      ótima, dá para \textbf{substituir uma pela outra} e treinar a política
      direto sobre os pares. {\scriptsize\color{rlCinza}(Rafailov et al., 2023.)}
    \end{ideiabox}
  \end{column}
  \end{columns}

  \vspace{-0.6em}
  \begin{center}
  \scalebox{0.72}{\begin{tikzpicture}[font=\scriptsize, node distance=7mm]
    \node[draw=rlCinza, thick, fill=rlFundo, rounded corners=2pt, minimum height=7mm,
          text width=20mm, align=center] (p) {preferências};
    \node[draw=rlTeal, thick, fill=rlTeal!10, rounded corners=2pt, minimum height=7mm,
          text width=20mm, align=center, right=11mm of p] (r) {recompensa $\rphi$};
    \node[draw=rlAzul, thick, fill=rlAzul!7, rounded corners=2pt, minimum height=7mm,
          text width=20mm, align=center, right=11mm of r] (pi) {política $\polth$};
    \draw[trans] (p) -- node[above, text=rlCinza] {entropia cruzada} (r);
    \draw[trans] (r) -- node[above, text=rlCinza] {PPO} (pi);
    \draw[trans destac] (p.south) .. controls +(0,-0.9) and +(0,-0.9) ..
      node[below, text=rlLaranja] {\textbf{DPO}: um único passo supervisionado} (pi.south);
  \end{tikzpicture}}
  \end{center}
\end{frame}

%% ------------------------------------------------------------
\begin{frame}{Passo 1: a solução em forma fechada}

  \begin{teobox}{Política ótima do objetivo com regularização KL}
    O máximo de
    $\;\E_{y\sim\pol}[r(x,y)] - \beta\,\KLd{\pol(\cdot\mid x)}{\pref(\cdot\mid x)}$
    é atingido por
    \[
      \pol_{r}(y\mid x) = \frac{1}{Z(x)}\,\pref(y\mid x)\,
        \exp\!\Bigl(\tfrac{1}{\beta} r(x,y)\Bigr),
      \qquad
      Z(x) = \sum_{y}\pref(y\mid x)\exp\!\Bigl(\tfrac{1}{\beta} r(x,y)\Bigr).
    \]
  \end{teobox}

  \pause
  \textbf{Prova em três linhas.} Divida por $\beta$ e complete o logaritmo:
  \begin{align*}
    \max_{\pol}\ \E_{y\sim\pol}\Bigl[\tfrac{1}{\beta}r - \log\tfrac{\pol}{\pref}\Bigr]
    &= \min_{\pol}\ \E_{y\sim\pol}\Bigl[\log\frac{\pol(y\mid x)}
        {\tfrac{1}{Z(x)}\pref(y\mid x)e^{r(x,y)/\beta}}\Bigr] - \log Z(x)\\
    &= \min_{\pol}\ \KLd{\pol(\cdot\mid x)}{\pol_r(\cdot\mid x)} - \log Z(x).
  \end{align*}
  Como $Z(x)$ não depende de $\pol$ e a divergência KL é mínima (e nula) apenas
  quando os argumentos coincidem, o mínimo é $\pol = \pol_r$. \hfill$\blacksquare$
\end{frame}

%% ------------------------------------------------------------
\begin{frame}{Passo 2: inverter a relação}

  A forma fechada é inútil como \emph{algoritmo}: $Z(x)$ soma sobre todas as
  respostas possíveis --- intratável. Mas ela pode ser \textbf{invertida}.

  \[
    \pol_{r}(y\mid x) = \frac{1}{Z(x)}\pref(y\mid x) e^{r(x,y)/\beta}
    \quad\Longleftrightarrow\quad
    \mdestaque{rlLaranja}{r(x,y) = \beta \log \frac{\pol_{r}(y\mid x)}{\pref(y\mid x)}
      + \beta\log Z(x)}
  \]

  \begin{ideiabox}
    Leitura: \emph{toda} função de recompensa pode ser escrita como uma razão
    logarítmica entre uma política e a referência. A razão é positiva se a
    política gosta da resposta \emph{mais} que o modelo de referência, e negativa
    se gosta menos.
  \end{ideiabox}

  \begin{itemize}\setlength{\itemsep}{0.25ex}
    \item O termo problemático $\beta\log Z(x)$ depende apenas de $x$.
    \item E o modelo de Bradley--Terry, felizmente, só enxerga \emph{diferenças}
      de recompensa para o mesmo $x$.
  \end{itemize}
\end{frame}

%% ------------------------------------------------------------
\begin{frame}{Passo 3: substituir --- e ver $Z$ desaparecer}

  Partindo de $\Prob(\yw \succ \yl \mid x) = \sig\bigl(r(x,\yw) - r(x,\yl)\bigr)$
  e substituindo a expressão anterior:
  \[
    r(x,\yw) - r(x,\yl)
      = \beta\log\frac{\polth(\yw\mid x)}{\pref(\yw\mid x)}
      - \beta\log\frac{\polth(\yl\mid x)}{\pref(\yl\mid x)}
      + \underbrace{\beta\log Z(x) - \beta\log Z(x)}_{\textstyle = \,0}
  \]

  \pause
  Levando à perda de entropia cruzada, obtém-se um objetivo que só envolve a
  política:
  \begin{teobox}{Perda do DPO}
    \[
      \mathcal{L}_{\mathrm{DPO}}(\polth;\pref) = -\,\E_{(x,\yw,\yl)\sim\Dset}
      \left[\log \sig\!\left(
        \beta\log\frac{\polth(\yw\mid x)}{\pref(\yw\mid x)}
        - \beta\log\frac{\polth(\yl\mid x)}{\pref(\yl\mid x)}
      \right)\right]
    \]
  \end{teobox}

  Uma perda \textbf{supervisionada}, sobre um conjunto \emph{fixo} de pares.
  Nenhuma amostragem, nenhum crítico, nenhum modelo de recompensa explícito.
\end{frame}

%% ------------------------------------------------------------
\begin{frame}{O que o gradiente do DPO faz}

  \[
    \scalebox{0.92}{$\displaystyle
    \gth \mathcal{L}_{\mathrm{DPO}}
      = -\beta\, \E_{(x,\yw,\yl)}\Bigl[\;
        \underbrace{\sig\bigl(\hat r_{\thv}(x,\yl) - \hat r_{\thv}(x,\yw)\bigr)}
          _{\textstyle \text{peso: quão errado o modelo está}}
        \Bigl(
          \underbrace{\gth\log\polth(\yw\mid x)}_{\text{aumentar}}
          - \underbrace{\gth\log\polth(\yl\mid x)}_{\text{diminuir}}
        \Bigr)\Bigr]$}
  \]
  com $\hat r_{\thv}(x,y) = \beta\log\frac{\polth(y\mid x)}{\pref(y\mid x)}$
  a \textbf{recompensa implícita} definida pela própria política.

  \begin{ideiabox}
    Estrutura familiar: é o gradiente de política ``suba o que deu certo, desça o
    que deu errado'' --- mas com um \textbf{peso adaptativo} que quase zera nos
    pares que o modelo já ordena corretamente e concentra o esforço nos que ele
    ainda erra. Esse peso é exatamente o que impede o colapso que ocorreria se
    treinássemos ingenuamente por máxima verossimilhança sobre $\yw$.
  \end{ideiabox}
\end{frame}

%% ------------------------------------------------------------
\begin{frame}{Teste sua compreensão}

  \begin{quizbox}{}
    O DPO nunca constrói um modelo de recompensa. Ainda assim, dizemos que ele
    tem uma ``recompensa implícita''.

    \smallskip
    (i) Onde ela está? (ii) O termo $\beta\log Z(x)$ cancelou --- isso significa
    que ele é irrelevante, ou que apenas não é \emph{necessário}? (iii) Que
    hipótese do RLHF o DPO herda inteira, sem enfraquecer?
    \paradiscussao
  \end{quizbox}

  \pause

  \begin{respbox}
    (i) Na razão $\beta\log(\polth/\pref)$: os pesos da política \emph{são} o
    modelo de recompensa reparametrizado. \;(ii) Ele existe e fixa a
    normalização; só não precisa ser conhecido, pois a perda depende apenas de
    diferenças com o mesmo $x$. \;(iii) O \textbf{modelo de Bradley--Terry}. Se
    ele falha, DPO e RLHF falham juntos.
  \end{respbox}
\end{frame}

%% ------------------------------------------------------------
\begin{frame}{O que se ganha e o que se perde}

  \begin{center}\small
  \begin{tabular}{@{}p{0.3\linewidth}p{0.31\linewidth}p{0.31\linewidth}@{}}
    \toprule
    & \textbf{RLHF (PPO)} & \textbf{DPO} \\
    \midrule
    Modelos em memória & 4 & 2 \\
    Dados & \emph{online} (regerados) & \emph{offline}, fixos \\
    Estabilidade & sensível a hiperparâmetros & perda supervisionada \\
    Rotular dados novos & possível a qualquer momento & exige nova coleta \\
    Compromisso recompensa/KL & explorado eficientemente & idem, com menos custo \\
    Explorar fora do conjunto & sim & \textbf{não} \\
    \bottomrule
  \end{tabular}
  \end{center}

  \begin{alertabox}
    A última linha é a limitação real. O DPO só pode reordenar o que já está no
    conjunto de dados; ele nunca descobre uma resposta melhor que ninguém
    escreveu. Daí as variantes \emph{iterativas}, que alternam geração,
    rotulagem e DPO.
  \end{alertabox}
\end{frame}

%% ------------------------------------------------------------
\begin{frame}{Evidência empírica e adoção}

  \begin{columns}[T]
  \begin{column}{0.49\textwidth}
    \begin{defbox}{Experimento controlado}
      \begin{enumerate}\setlength{\itemsep}{0ex}
        \item Gerar resenhas positivas de filmes com GPT-2 XL.
        \item Usar um classificador de sentimento como recompensa ``de ouro''.
        \item Construir preferências sintéticas a partir dela.
        \item Otimizar com PPO e com DPO.
      \end{enumerate}
      Resultado: para o \emph{mesmo} desvio KL, o DPO obtém recompensa igual ou
      maior --- fronteira de compromisso melhor.
    \end{defbox}
  \end{column}
  \begin{column}{0.49\textwidth}
    \begin{itemize}\setlength{\itemsep}{0.3ex}
      \item Com recompensa de ouro conhecida, a comparação é limpa: não há
        dúvida sobre a qualidade do modelo de recompensa.
      \item Em escala, o DPO e suas variantes passaram a ser padrão em famílias
        abertas de modelos (Zephyr, Mistral, Llama 3).
      \item Extensões próximas: \textbf{CPL}, que leva a mesma ideia para
        preferências sobre \emph{trajetórias} de controle, fora do domínio de
        linguagem.
    \end{itemize}
  \end{column}
  \end{columns}
\end{frame}

%% ------------------------------------------------------------
\begin{frame}{O mapa depois do DPO}

  \begin{center}
  \scalebox{0.83}{\begin{tikzpicture}[font=\scriptsize, node distance=6mm]
    \node[draw=rlAzul, very thick, fill=rlAzul!7, rounded corners=3pt,
          text width=26mm, align=center, minimum height=11mm] (bt)
      {\textbf{Bradley--Terry}\\ preferências $\to$ recompensa};
    \node[draw=rlTeal, thick, fill=rlTeal!10, rounded corners=3pt,
          text width=26mm, align=center, minimum height=11mm, above right=3mm and 16mm of bt] (rlhf)
      {\textbf{RLHF}\\ recompensa explícita $+$ PPO};
    \node[draw=rlLaranja, very thick, fill=rlLaranja!10, rounded corners=3pt,
          text width=26mm, align=center, minimum height=11mm, below right=3mm and 16mm of bt] (dpo)
      {\textbf{DPO}\\ recompensa implícita, perda direta};
    \node[draw=rlAzulClaro, thick, fill=rlAzulClaro!8, rounded corners=3pt,
          text width=30mm, align=center, minimum height=11mm, right=16mm of rlhf] (grpo)
      {\textbf{GRPO / RLVR}\\ recompensa \emph{verificável}: dispensa o modelo};
    \node[draw=rlAzulClaro, thick, fill=rlAzulClaro!8, rounded corners=3pt,
          text width=30mm, align=center, minimum height=11mm, right=16mm of dpo] (var)
      {\textbf{IPO, KTO, ORPO}\\ afrouxam hipóteses de Bradley--Terry};
    \draw[trans destac] (bt) -- (rlhf);
    \draw[trans destac] (bt) -- (dpo);
    \draw[trans] (rlhf) -- (grpo);
    \draw[trans] (dpo) -- (var);
  \end{tikzpicture}}
  \end{center}

  \begin{ideiabox}
    O ramo de cima move o estado da arte em raciocínio: quando a correção pode
    ser \emph{verificada} (uma prova, um teste unitário, uma resposta numérica),
    a recompensa volta a ser dada --- e os problemas de sobre-otimização desta
    aula desaparecem.
  \end{ideiabox}
\end{frame}

%% ============================================================
\section{Balanço}
%% ============================================================

%% ------------------------------------------------------------
\begin{frame}{O que você deve levar desta aula}

  \begin{defbox}[fontupper=\small]{Capacidades esperadas}
    \begin{itemize}\setlength{\itemsep}{0ex}
      \item Explicar a degradação $O(\varepsilon T^2)$ da clonagem, e por que
        mais dados não corrigem o expoente.
      \item Descrever o DAGGER, sua garantia e por que ele não escala.
      \item Argumentar que a recompensa do RL inverso não é única; listar as
        três invariâncias.
      \item Escrever Bradley--Terry e a perda correspondente, para itens e para
        trajetórias.
      \item Descrever as três etapas do RLHF e o papel da penalidade KL.
      \item \textbf{Derivar} o DPO e interpretar seu gradiente.
    \end{itemize}
  \end{defbox}

  \begin{ideiabox}
    A moral: quando não há placar, alguém precisa fornecê-lo. A aula toda foi
    sobre \emph{quem} fornece, \emph{a que custo} e \emph{quanto se pode
    confiar} no que veio.
  \end{ideiabox}
\end{frame}

%% ------------------------------------------------------------
\begin{frame}{Para ir além}

  \begin{columns}[T]
  \begin{column}{0.49\textwidth}
    \textbf{Artigos fundamentais}
    \begin{itemize}\setlength{\itemsep}{0ex}
      \item Ross, Gordon \& Bagnell (2011) --- DAGGER.
      \item Ng \& Russell (2000) --- RL inverso.
      \item Ziebart et al. (2008) --- MaxEnt IRL.
      \item Christiano et al. (2017) --- preferências humanas.
      \item Ouyang et al. (2022) --- InstructGPT.
      \item Rafailov et al. (2023) --- DPO.
    \end{itemize}
  \end{column}
  \begin{column}{0.49\textwidth}
    \textbf{Contexto mais amplo}
    \begin{itemize}\setlength{\itemsep}{0ex}
      \item Decidir a partir de preferências humanas fica na interseção de
        \textbf{escolha social}, \textbf{economia computacional} e IA, com
        literatura anterior ao aprendizado de máquina.
      \item Teoremas de impossibilidade (Arrow) afetam o que significa
        ``alinhar'' um modelo a um \emph{grupo} com preferências conflitantes.
    \end{itemize}
  \end{column}
  \end{columns}

  \begin{alertabox}
    Uma pergunta que a técnica não responde: \emph{de quem} são as preferências
    usadas para rotular os pares. É escolha de projeto, não detalhe de
    implementação.
  \end{alertabox}
\end{frame}

\end{document}
